%! Comparing Studies and Choosing a Method — Unit 1, Lesson 1.7 — Lesson Plan
% PILOT SHAPE (2026-09-06), notes in the MAIN IDEAS / NOTES table shape (2026-09-12).
% GRADUAL-RELEASE plan (warm-up / guided notes and practice / spoken debrief / close and START
% THE HOMEWORK). The guided notes carry the release ROW BY ROW in one table: I do / we do in
% rows 1-4 (problems 1-15) -> we do the Guided Practice row (problems 16-19). THERE IS NO
% INDEPENDENT PRACTICE SET IN THE NOTES — the homework is the individual practice, and the
% period ends with students starting it. No group activity, no tier boxes, no exit ticket:
% the old activity's crux became the Guided Practice row and the old exit ticket's conceptual
% item became the debrief's cold check. Teacher-facing.
\documentclass[10pt]{article}
\usepackage{saar-article}
\usepackage{saar-boxes}
\usepackage{graphicx}   % -article does NOT load graphicx
\graphicspath{{images/}}
\newcommand{\TallMath}[1]{$\displaystyle #1\rule[-1.4em]{0pt}{3.2em}$}
\newcommand{\UnitNumberName}{Unit 1: Asking Questions with Data: Study Design \quad}
\newcommand{\LessonNumberName}{Lesson 1.7: Comparing Studies and Choosing a Method}

\begin{document}
\begin{center}
    {\Large\bfseries \CourseName \\
    {\small\bfseries \UnitNumberName \LessonNumberName}}
\end{center}

\begin{tcolorbox}[colback=frost, colframe=cerulean, boxrule=0.9pt, arc=2mm,
    left=2mm, right=2mm, top=1.5mm, bottom=1.5mm]
    \textbf{Primary Objective:} Students will put an observational study and a controlled
    experiment side by side, write the strongest sentence each design earns, decide whether an
    experiment can be run at all --- or whether assigning the treatment is unethical,
    impossible, or impractical --- plan a well-designed experiment through the four stages of
    the statistical cycle, and choose the collection method a question needs while keeping the
    \emph{method} separate from the \emph{design}.
    \\[2pt]
    \textbf{Standards:} PS.DC.3c, PS.DC.3d, PS.DC.3e \quad$\cdot$\quad PS.DC.1b carried forward
    in the conclusion items \quad$\cdot$\quad reaches back to 1.5 (sampling) and 1.6
    (experiments)
    \\[2pt]
    \textbf{Lesson model:} \emph{Gradual release --- I do, we do, you do --- all of it inside
    the guided notes.} There is no group activity and no exit ticket. The notes name the
    vocabulary as they teach it and deliver the instruction \textbf{row by row in one Main
    Ideas / Notes table} (four instruction rows, problems 1--15: you read the definition, mark
    up the display, and work the first problem of each grid; the class works the rest); the
    Guided Practice row (problems 16--19, \emph{Harbor Point}) is worked together with students
    holding the pen; \textbf{the homework is the individual practice}, and students start it in
    class while you circulate. The whole-class debrief is \emph{spoken}.
    \\[2pt]
    \textbf{Note on scope:} the four principles of a good experiment (comparison, randomization,
    replication, control) are \emph{assumed} from Lesson 1.6, not re-taught --- they are named
    once, inside stage 2 of the cycle. Nothing today requires a new computation beyond a percent
    and a gap.
\end{tcolorbox}

\renewcommand{\arraystretch}{1.4}
\boxguard
\begin{skillbox}[Learning Targets \& Key Understandings]{goldbox}
\begin{tabularx}{\linewidth}{@{}X|X@{}}
\textbf{Learning targets (student language)} & \textbf{Key understandings (the ``why'')} \\[2pt]
\hline
\vspace{-6pt}
\begin{itemize}[leftmargin=1.1em, nosep, after=\vspace{2pt}]
  \item I can compare an \textbf{observational study} with a \textbf{controlled experiment}
        side by side, and write the strongest sentence each design earns.
        \hfill \textit{PS.DC.3c}
  \item I can decide whether an experiment can be run at all --- or name it
        \textbf{unethical}, \textbf{impossible}, or \textbf{impractical} --- and say what to
        run instead. \hfill \textit{PS.DC.3c}
  \item I can plan a well-designed experiment through the four stages of the
        \textbf{statistical cycle}. \hfill \textit{PS.DC.3d}
  \item I can choose the \textbf{collection method} a question needs --- survey, interview,
        focus group, observation, content analysis. \hfill \textit{PS.DC.3e}
  \item I can keep the \emph{method} separate from the \emph{design}, and say what a gap does
        and does not prove. \hfill \textit{PS.DC.3c, 3e; 1b}
\end{itemize}
&
\vspace{-6pt}
\begin{itemize}[leftmargin=1.1em, nosep, after=\vspace{2pt}]
  \item \emph{One question produces the whole comparison table: who decided which group each
        individual is in?} Every other row --- treatment, groups alike, the verb --- follows
        from that answer.
  \item \textbf{Target misconception:} that \emph{observation} the collection method makes a
        study an \emph{observational study}. It does not. An experiment can collect its data by
        observation: the counselor watched who came, but a generator built her groups, so the
        design is an experiment (problem~15, HW~6(a)).
  \item \textbf{Second misconception:} that the biggest gap is the best evidence. Millbrook's
        $48$-point gap is the biggest number in the unit and its weakest evidence --- the card
        holders sorted themselves (HW~6(b)).
  \item ``We could not run an experiment'' is a reason for a \textbf{weaker} verb, never an
        excuse for a stronger one (problem~9).
  \item An experiment is the stronger design, so the real question is never ``which is better''
        but \textbf{which one can I actually run} --- can you hand this out to a person?
\end{itemize}
\end{tabularx}
\end{skillbox}

\boxguard[24]
\begin{skillbox}[{Vocabulary, Concepts, \& Theorems --- taught directly in the guided notes}]{greenbox}
\small
\renewcommand{\arraystretch}{1.35}
\begin{tabularx}{\linewidth}{@{} >{\bfseries}p{3.1cm} X @{}}
Controlled experiment & A study in which the researcher assigns the units to groups by chance
       and imposes a treatment on at least one group. It can support a cause-and-effect
       conclusion. \\
Observational study & A study in which the researcher measures or asks without assigning anyone
       to a group --- the groups already exist. It can report an association only. \\
Collection method & How the data arrive: survey, interview, focus group, observation, or
       content analysis. Chosen by what the question needs, not by the design. \\
Statistical cycle & Formulate $\rightarrow$ collect $\rightarrow$ represent $\rightarrow$
       analyze \& report. The same four stages plan either study type; only stage 2 changes. \\
Survey & The same fixed questions given to every individual in a chosen sample. Best for counts,
       percents, and attitudes across many people. \\
Interview & A one-on-one conversation built on planned questions. Best for depth --- the reasons
       behind an answer, in the person's own words. \\
Focus group & A small-group discussion led by a facilitator, where members build on each other's
       answers. \\
Observation & Recording what people actually do. Best when behavior and self-report may
       disagree. \emph{A method, not a design.} \\
Content analysis & Studying records or documents that already exist --- borrowing histories,
       transcripts, past scores. Best for change over time. \\
\end{tabularx}
\end{skillbox}

% fixedskillbox never splits, so the phase table stays intact on one page.
% THE PHASE TABLE IS A CONTRACT: 5 / 35 / 10 / 10 — the I do is ~20 of the 35 and the Guided
% Practice ~15; the debrief is spoken; the homework start is the second formative read.
\begin{fixedskillbox}[Lesson at a Glance --- \MeetingLength]{frost}
\renewcommand{\arraystretch}{1.25}
\begin{tabularx}{\linewidth}{@{}l c X X@{}}
\textbf{Phase} & \textbf{Min} & \textbf{Students} & \textbf{Teacher} \\
\hline
Warm-Up & 5 & Two Riverbend percents, then the two studies side by side --- silent & Debrief
  item~3 aloud; leave \emph{$60$ and $90$} against \emph{$60$ and $60$} on the board all
  period \\
Guided Notes \& Practice & 35 & Fill the vocabulary box as each term is named; work rows 1--4
  (problems 1--15) with the pen in their hand; then the Guided Practice row, \emph{Harbor
  Point} (16--19), together & \textbf{I do / we do, row by row}: read the definition, mark up
  the display, work problems 2, 6, 10, 12; the class works the rest (20) $\to$ \textbf{we do}
  Guided Practice 16--19 (15) \\
Debrief & 10 & Pens down, papers up --- answer aloud, nothing to fill in & Put the counselor
  back on the board and take the \emph{``she just watched''} reflex, then the
  \emph{``unethical''} reflex from problem~8 \\
Close \& Assign & 10 & \textbf{Start the homework in class}, alone and silent & Announce the
  homework and its form (packet or Desmos), circulate for the second formative read, preview
  Lesson~1.8 \\
\end{tabularx}
\end{fixedskillbox}

\begin{fixedskillbox}[Warm-Up --- Activate Prior Knowledge (5 min)]{frost}
\begin{tabularx}{\linewidth}{@{}X X@{}}
\begin{minipage}[t]{\linewidth}\vspace{0pt}
    \textbf{The three items and what each seeds}
    \begin{itemize}[leftmargin=1.1em, itemsep=1pt, topsep=2pt]
      \item \textbf{1.} $39$ of $60$ as a percent ($65\%$). \textbf{Spirals} to Lesson 1.5 ---
            and it is the number \emph{problem~2} labels with the verb \emph{is associated
            with}.
      \item \textbf{2.} $72$ of $90$ ($80\%$). \textbf{Seeds:} the ``harmful'' result ---
            \emph{problem~4} asks which of the two studies to believe once both are on the
            page.
      \item \textbf{3.} The two studies side by side, no vocabulary: the two gaps ($15$ and
            $20$) and the group sizes. \textbf{Seeds:} row~1 --- and its last question is
            \emph{problem~5} in advance.
    \end{itemize}
\end{minipage}
&
\begin{minipage}[t]{\linewidth}\vspace{0pt}
    \textbf{Running it}
    \begin{itemize}[leftmargin=1.1em, itemsep=1pt, topsep=2pt]
      \item \textbf{Debrief item~3 aloud} and write the two group sizes on the board ---
            \emph{$60$ and $90$} against \emph{$60$ and $60$}. Leave it up all period; row~1
            hangs both design names on that line and problem~5 comes straight back to it.
      \item Item~3's last question is the tell. A student who answers \emph{the students
            themselves} has already drawn today's distinction without the words and row~1 can
            move fast; a student who writes \emph{the counselor} needs row~1 slowly.
      \item Items 1--2 are Lesson 1.5's arithmetic --- say so, so nobody thinks they are being
            retested. Do not let them expand.
    \end{itemize}
\end{minipage}
\end{tabularx}
\end{fixedskillbox}

\boxguard
\begin{skillbox}[Hook --- before anyone sits down]{frost}
The hook is on \textbf{page 1 of the notes}, under the vocabulary box, and on the board:
\textit{``This quarter a random number generator split $120$ volunteers into two groups of
$60$. To collect her data the counselor handed out no questionnaire --- she sat in the Study
Center three afternoons a week and \emph{watched}. A classmate says: she just watched, so this
quarter's study was an observational study, the same as last quarter's.''} \textbf{Ask: is the
classmate right?} Students circle \emph{agree} or \emph{disagree} and write one line of reason
\emph{before} anyone speaks; then take the hand vote and write the split on the board. Most
rooms agree with the classmate --- watching sounds like observing.
\textbf{Do not settle it.} Say only that the answer is \textbf{problem~15}, in the
\emph{The Method} row, and that they will vote again there. The vote is what makes the crux
land: students have to have committed to \emph{watching makes it observational} before the notes
take it away.
\end{skillbox}

\boxguard
\begin{skillbox}[Guided Notes \& Practice --- I do, then we do (35 min)]{frost}
\textbf{This is the whole lesson.} There is no group activity, and there is no independent
practice set on this page: \textbf{the homework is the individual practice}, started in class.
The notes are \textbf{one Main Ideas / Notes table}: four instruction rows (problems 1--15),
then the Guided Practice row (16--19). \textbf{The rhythm in every row is the same}: read the
printed sentences aloud, mark up the display, work the first problem of the grid yourself, then
hand the rest to the room. The vocabulary box on page~1 is filled \emph{as each term is named},
never in advance. The only blanks are the six cells of the side-by-side table (problem~1);
every other answer is written in a problem's own space.
\begin{multicols}{2}
\textbf{I do / we do, row by row (about 20 min).}

\emph{Row 1, Two Designs (problems 1--5).} Open on the line already on the board: $60$ and $90$
against $60$ and $60$. Read the two definitions, then name the two panels of the display ---
left an \emph{observational study}, right a \emph{controlled experiment}. Take problem~1's table
from the class row by row; the first row is given because \textbf{it produces every other row}.
Work problem~2 yourself (\emph{is associated with}); the class does 3. Problems 4 and~5 are the
payoff --- which result to believe, and how the group sizes gave the designs away. Six minutes.

\emph{Row 2, Experiment --- when you cannot (problems 6--9).} Give the three reasons and the
test question, then walk the flowchart. Work problem~6 yourself; 7 goes to the room.
\textbf{Problem~8 is the trap, and it is worth a hand count before anyone writes}: the
$45$-minute bus ride. Most of the room will accept the classmate's \emph{``unethical.''} It is
not --- it is \textbf{impossible}: nobody can assign where a family lives. Problem~9 lands the
rule that an unavailable experiment weakens the verb. Six minutes.

\emph{Row 3, The Cycle (problems 10--11).} Fast --- four minutes. The four stages are Lesson
1.5's; only stage~2 changes, from \emph{sample} to \emph{assign}. Work problem~10 yourself,
hand 11 to the room, and move.

\emph{Row 4, The Method (problems 12--15) --- the crux.} Read the five method cards, land the
two-question box (\emph{who built the groups?} $\to$ design; \emph{how were the answers
recorded?} $\to$ method), and take the \emph{In your own words} line before the grid. Work
problem~12 yourself, let the class do 13 and~14 --- and then \textbf{re-take the hook vote at
problem~15 before you say anything}. She watched (method: observation) inside a study a
generator built (design: experiment). Four minutes on 12--14, four on 15.

\columnbreak
\textbf{We do (about 15 min).} The Guided Practice row, \emph{Harbor Point} (problems
16--19) --- the same moves the homework will ask alone, on a pool instead of a school.
\textbf{Students hold the pen; you hold the questions.} Problem~16 goes on them cold: label the
1.5 survey and the 1.6 experiment and say who built the groups in each. Problem~17 is warm-up
arithmetic with a verb attached ($42/60 = 70\%$ against $30/60 = 50\%$, gap $20$).
\textbf{Problem~18 is the cannot-assign case transferred} --- arthritis is something a member
already \emph{has}. \textbf{Problem~19 is the crux transferred}: a survey and interviews inside
an experiment, and neither changes the design.

Circulate during Guided Practice and demand the reason, not the word:
\begin{itemize}[leftmargin=1.1em, nosep]
  \item \textbf{Problem 16:} ``Who decided which group each member was in?'' Nothing else sorts
        the two studies.
  \item \textbf{Problem 17:} ``You wrote a $20$-point gap. Which verb did the design buy you?''
  \item \textbf{Problem 18:} ``Show me how you would hand arthritis out to a volunteer.''
  \item \textbf{Problem 19:} ``Who built the groups? Then what is it, no matter how the answers
        arrived?''
\end{itemize}
While circulating, pick the papers to put up in the debrief --- one problem~19 that kept the
design and one that switched it to observational. \textbf{Your formative read comes twice}:
here, and again in the last ten minutes as students start the homework.
\end{multicols}
\end{skillbox}

\boxguard[24]
\begin{skillbox}[Debrief --- whole class, spoken (10 min)]{frost}
Pens down, papers up. \textbf{This debrief is spoken} --- there is no box at the end of the
notes to fill in, so it lives entirely on the board and in the room.
\begin{multicols}{2}
\begin{enumerate}[leftmargin=1.3em, nosep, itemsep=3pt]
  \item \textbf{Put the counselor back up} --- she watched, a generator built her groups --- and
        make a student say the sentence: \emph{observation is how the data arrived; the
        generator is what makes it an experiment}. Then state it as PS.DC.3e: the collection
        method is chosen by the question, and it never decides the design.
  \item Put up the two problem~19s you picked. The difference between them is not arithmetic ---
        it is whether the student still sorts studies by how the data were recorded.
  \item Put problem~8 back up and take the \emph{unethical} reflex: \emph{``Who could hand out a
        $45$-minute bus ride? Then what is the reason?''} (Impossible, not unethical.)
  \item Take the verb last, from problems 2--3: the same Study Center, two studies, two
        sentences --- \emph{is associated with} and \emph{caused}. What bought the second one?
\end{enumerate}
\columnbreak
\textbf{Check for understanding --- ask it cold, whole class.} \emph{``Bayside Middle School has
two studies. Study A: the counselor compared the $150$ students in band with the $600$ who are
not --- $90$ of the band students had a B average. Study B: $80$ sixth-graders volunteered and a
generator sent $40$ to a $10$-minute silent reading block and $40$ to their usual homeroom;
$70\%$ improved against $40\%$. Which study may say the program \emph{caused} the difference,
and what percent of the band students had a B average?''} Hands down, thirty seconds, then
cold-call three students. Correct answer: \textbf{Study B} --- a generator built its groups;
Study A's groups were already there. And $90/150 = \textbf{60\%}$. \emph{This replaces the exit
ticket.}

\textbf{Formative read.} Sort the three answers: \textbf{(a)} picked B and said \emph{a
generator built the groups}; \textbf{(b)} picked B with no reason; \textbf{(c)} picked A, or
sorted by how the data were collected. Pile (b) is ready and needs the language. \textbf{If pile
(c) runs past a third of the room, open Lesson~1.8 by re-running the \emph{The Method} row
(problems 12--15)}, and say so plainly.

\textbf{If the debrief runs short}, cut items 3 and~4 --- item~1 is the one that cannot be cut.
\textbf{Do not let the debrief eat the last ten minutes} --- the homework start is not optional
padding, it is your second formative read.
\end{multicols}
\end{skillbox}

\boxguard
\begin{skillbox}[Homework --- scored, started in class, due after two study halls]{goldbox}
Six exercises on the \textbf{Millbrook Public Library} ($1{,}200$ card holders; a \emph{third}
study --- $300$ volunteers, $150$/$150$, $96$ against $60$ reading at least $10$ books):
\textbf{1} sort the two earlier studies by design and say who built the groups; \textbf{2}
compute $64\%$ and $40\%$, the $24$-point gap, and write the sentence the design earns;
\textbf{3} the book-club study through the four stages of the cycle (\emph{the spiral to Lesson
1.5's cycle table}); \textbf{4} two new board questions --- the late fee against being read to
as a toddler: the \textbf{contrast pair}, one assignable and one not; \textbf{5} all five
collection methods matched to five jobs; \textbf{6} \emph{the crux head-on} --- (a) she
collected the data by sitting in and watching, so a classmate calls it an observational study,
and (b) the $48$-point gap, the biggest number in the unit and the weakest evidence.

\textbf{Start it in the last ten minutes of the period, in class and alone.} \textbf{Scoring:}
12 points --- 2 per item. \textbf{Item~6(a) is where the misconception shows}: a student who
agrees with the classmate has learned the five methods and not the rule. \textbf{Item~6(b)
carries the second check}, and item~4 row~(ii) the trap reason (\emph{impossible}, not
\emph{unethical}).

\textbf{Packet or Desmos --- decided here.} The packet pages are authored and are the default.
Desmos Classroom's study-design coverage is thin --- nothing in it asks a student to sort a
design from a method or to rank three gaps by strength of evidence --- so \textbf{1.7 is a
packet night}. Say so aloud at Close \& Assign.

\textbf{Preview of Lesson~1.8.} Every design decision students judged today, they make
themselves next time, on a question the class chooses.

\textbf{Connections \& Big Ideas.} \textit{Unit 1 core idea:} a conclusion is only as
trustworthy as the study that produced it. \textit{Standards covered:} PS.DC.3c, PS.DC.3d,
PS.DC.3e (with PS.DC.1b carried forward in the conclusion items). \textit{Spirals forward to:}
the full-cycle project (1.8, AFDA.DA.2d, 2h--j) and, directly, inference (7.1--7.2).
\textit{Reaches back to:} Lesson 1.5's sampling and cycle, Lesson 1.6's four principles, and
Algebra~1 percent reasoning.
\end{skillbox}

\boxguard
\begin{skillbox}[Watch For (while circulating)]{redbox}
\begin{itemize}[leftmargin=1.2em, nosep]
  \item \textbf{``She watched, so it is an observational study''} (problem~15, problem~19,
        HW~6(a) --- the target misconception). Probe: \emph{``Who built the groups? Does how
        the answers were recorded change that?''}
  \item \textbf{``Unethical'' used for anything you cannot assign} (problem~8, HW~4 row (ii)).
        Probe: \emph{``Who is harmed by a long bus ride being studied? Now --- could you hand
        one out?''}
  \item \textbf{The biggest gap is the best evidence} (problem~4, HW~6(b)). Probe: \emph{``Who
        put the card holders into those two groups?''}
  \item \textbf{``No experiment was possible, so we may still say caused''} (problem~9). Probe:
        \emph{``Does not being able to run it make the groups alike?''}
  \item \textbf{Survey chosen for every job} (problems 12--14, HW~5). Probe: \emph{``Will a
        written question tell you what people actually did?''}
  \item \textbf{The cycle recited, not planned} (problems 10--11, HW~3): \emph{``collect the
        data''} as a stage answer. Probe: \emph{``Who are the units, and what does the
        generator do?''}
\end{itemize}
Cold-call prompts: \emph{``Name a question about this school that nobody could ever answer with
an experiment.''} \quad \emph{``Give me an experiment whose data arrive by observation.''}
\end{skillbox}

\boxguard
\begin{skillbox}[Close \& Assign (10 min)]{goldbox}
\textbf{Students start the homework here, in class, alone and silent --- this is the individual
practice, and the first minutes of it are your best formative read.} Hand the launch first ---
six exercises on paper, scored out of 12, due at the start of the first class after two study
halls --- and point out that the \emph{Keep in Mind} box on the cover is the page to flip back
to. Circulate: \textbf{item~6 is the column that tells you whether \emph{The Method} row
landed}, and it is on page~2, so put students on 1, 2, and~6 first and let 3, 4, and~5 go home.
Sort what you see into three piles as you walk --- named the method and kept the design
(ready); right verdict, no reason (ready, needs the language); called the book-club study
observational (the misconception survived) --- and open Lesson~1.8 by re-running problems 12--15
if that third pile ran past a third of the room. Then close on what changed: \emph{``You walked
in able to tell a study that assigns from a study that watches. You walk out knowing that how
the data arrive --- a survey, an interview, somebody sitting in the corner watching --- never
decides which one it is, and that the biggest gap in this unit is the one you should trust
least.''} \textbf{Preview:} Lesson~1.8, \emph{You Run the Study} --- today you judged other
people's designs; next time the question is yours, and every decision on this page is one you
have to make.
\end{skillbox}

% ── Teacher notes — one per component, in packet order (three) ──────────────
% Teacher-only prose lives HERE, never in a _key. No note for the debrief or the homework
% start — both are fully specified in their own boxes above.
\begin{teachernote}[Warm-Up]
Five minutes, silent, timed. Items 1 and~2 are Lesson 1.5's arithmetic on purpose --- say so, so
nobody thinks they are being retested. \textbf{Item~3 is the one to read while collecting}:
students who write the two gaps ($15$ and $20$) and then answer \emph{the students themselves}
have already made today's distinction without the vocabulary, and row~1 can move fast. Students
who name the counselor as the one who put students into groups last quarter have the two studies
fused and will need row~1 slowly. Write the two group sizes on the board --- \emph{$60$ and
$90$} against \emph{$60$ and $60$} --- and \textbf{leave them up all period}; row~1 hangs both
design names on that line, problem~5 comes straight back to it, and you circle it again in the
debrief. Do \emph{not} supply the words \emph{observational study} and \emph{experiment} here.
\end{teachernote}

\begin{teachernote}[Guided Notes \& Practice]
\textbf{This one document is the lesson --- pace it 20 / 15, then 10 for the debrief, then 10 for
the homework start.} It is one table, and every row runs the same way: you read the printed
sentences, mark up the display, and work the first problem of the grid (2, 6, 10, 12); the class
works the rest with the pen in their hand. \textbf{Row~3, The Cycle, is the row to keep short} ---
four minutes. It is Lesson 1.5's table with one word changed, and if you teach it as new
material you will not reach problem~15 with time to argue.

\textbf{Take the hand count on problem~8 before anyone writes.} If students write first, the trap
is gone; if they commit to \emph{unethical} and \emph{then} get asked who is harmed, the
distinction arrives as the answer to a question they already have. The reason is
\emph{impossible} --- a treatment nobody can hand out --- and it is the same shape as HW~4
row~(ii) and Guided Practice problem~18.

\textbf{Spend the time in row~4, \emph{The Method}}, which carries the misconception. Read the
five cards quickly --- they are a reference table, not five lessons --- and put the weight on the
two-question box and on problem~15. \textbf{Re-take the hook vote at problem~15 before you say
anything}: agree or disagree. Let the room argue. Then separate the two questions out loud:
\emph{who built the groups?} (a generator --- experiment) and \emph{how were the answers
recorded?} (she watched --- observation). Do not resolve it by repeating the definition; ask the
two questions in order and let the answers disagree in front of them.

\textbf{Guided Practice (16--19) is where the pen changes hands.} \emph{Harbor Point} runs the
four moves the homework will ask alone --- label the designs, compute and attach the verb, the
cannot-assign case, the method-inside-an-experiment case --- on a pool instead of a school. If
you only get through 16, 17, and~19, that is the right trade; 18 can be said aloud.

\textbf{There is no independent practice set on this page, and that is deliberate --- the
homework is the individual practice.} Guided Practice is the last row; the period ends with
students starting the homework in front of you, which is where your second formative read comes
from. If the \emph{``she just watched''} pile is more than a third of the room, \textbf{open
Lesson~1.8 by re-running problems 12--15}, and say so plainly.
\end{teachernote}

\begin{teachernote}[Homework]
\textbf{Scored, 12 points (2 per item), due at the start of the first class after two study
halls.} The ten minutes at the end of the period are a \emph{start}, not a completion: put
students on 1, 2, and~6 while you can still catch a wrong start, and let 3, 4, and~5 go home.
Item~6 is the one to read first when scoring. In 6(a), any student who agrees with the classmate
is still sorting studies by how the data were recorded; in 6(b), the expected wrong answer is
that the $48$-point gap is the strongest evidence because it is the biggest number --- the
card holders signed themselves up, so the gap measures who signed up. Item~4 row~(ii) is the
second diagnostic: \emph{impossible}, not \emph{unethical}. Score items 1--5 for correctness and
item~6 for the \emph{reason}, not the verdict. \textbf{When you score the page, sort item~6(a)
into three piles} --- named the method and kept the design (ready); right verdict, no reason
(ready, needs the language); called it an observational study (the misconception survived). If
more than a third land in the last pile, reopen the counselor in the Lesson~1.8 warm-up debrief
rather than letting it ride.

\textbf{Packet is the call for 1.7.} Desmos carries nothing that asks a student to separate a
method from a design; do not swap it in. Millbrook's three studies are the last time this unit
hands students a finished design --- in 1.8 they build one.
\end{teachernote}
\end{document}
